\begin{proof}
\begin{enumerate}[label=\arabic*.]
\item Let $L=H^1(C_\lambda,\mathbb Z)$ with its integral cup pairing,
and write $T_1,T_2$ for the finite generators of
\cref{maps:src:F.3.5}. Put $A=C_f$ and $B=C_g$.
The rational intertwiner proved there is an actual map
\[
 q:L_{\mathbb Q}\longrightarrow\mathbb Q^4,
 \qquad qT_1q^{-1}=A^{-1},\quad qT_2q^{-1}=B.
\]
Its relation to the Levelt presentation is
$(h_\infty,h_1,h_0)=(T_1^{-1},T_1T_2,T_2^{-1})$;
the ordered product is one. The known rigidity input and its explicit
rational change of basis are recorded in \cref{maps:src:F.3.5}.
The nodal peripheral is $c=(T_1T_2)^{-1}$. Its primitive
Picard--Lefschetz class $\delta\in L$ and coefficient one are the
geometric node input proved there: the nonseparating vanishing circle
has a transverse integral cycle meeting it once, and the annulus
$xy=t$ turns once. Thus $T_1T_2$ acts by
$v\mapsto v+\epsilon\delta\langle\delta,v\rangle$,
where $\epsilon\in\{1,-1\}$ records the fixed inverse and pairing
conventions.

\item A $B$-invariant alternating form has upper entries
$(u,v,u,u,v,u)$. Invariance under $A$, evaluated on the first and
fourth basis vectors, gives $v=-u$. The common invariant forms are
therefore exactly the multiples of
\[
 J=\begin{pmatrix}
 0&1&-1&1\\-1&0&1&-1\\1&-1&0&1\\-1&1&-1&0
 \end{pmatrix},\qquad
 d=\begin{pmatrix}1\\1\\1\\0\end{pmatrix}.
\]
The Pfaffian of $J$ is $1$, and direct multiplication gives
\[
 A^tJA=B^tJB=J,\qquad
 A^{-1}B-I=d e_4^t,\qquad d^tJ=e_4^t.
\]
Here $e_4$ is the fourth standard vector. Write the transported cup
form as $sJ$ and $q(\delta)=k d$, with $s,k\in\mathbb Q^*$.
Comparison with the actual Picard--Lefschetz formula gives
$\epsilon s k^2=1$.

\item The integral classes
$\delta_i=T_1^{-i}\delta$, $0\leq i\leq3$, have rational images
equal to the columns of $kD$, where
\[
 D=(d,Ad,A^2d,A^3d)=
 \begin{pmatrix}1&0&-1&0\\1&1&-1&-1\\1&1&0&-1\\0&1&0&0\end{pmatrix},
 \qquad \det D=1.
\]
Their actual pairing is $\epsilon G$, with
\[
 G=D^tJD=
 \begin{pmatrix}0&1&0&0\\-1&0&1&0\\0&-1&0&1\\0&0&-1&0\end{pmatrix},
 \qquad \det G=1.
\]
These four classes are a basis of $L_{\mathbb Q}$. For any $v\in L$,
write $v=\sum_i x_i\delta_i$ with $x_i\in\mathbb Q$.
The vector $(\langle\delta_i,v\rangle)_i=\epsilon Gx$ is integral.
Since $G^{-1}$ is integral, so is $x$. Thus the four classes are an
integral basis of $L$.

\item The required integral marking is the precise map
$\mu=k^{-1}q:L\longrightarrow\mathbb Z^4$.
Indeed $q(L)=kD\mathbb Z^4=k\mathbb Z^4$. Its inverse is specified
on standard vectors by
\[
 \mu^{-1}(e_j)=\sum_{i=0}^3(D^{-1})_{i+1,j}\delta_i,
 \qquad
 D^{-1}=\begin{pmatrix}1&-1&1&0\\0&0&0&1\\0&-1&1&0\\1&-1&0&1\end{pmatrix}.
\]
This supplies both integral arrows and retains the cup form
$\epsilon J$. In this marking the finite pair is $(A^{-1},B)$.

\item The exact simultaneous companion reversal is
\[
 W=\begin{pmatrix}0&0&0&1\\0&0&1&0\\0&1&0&0\\1&0&0&0\end{pmatrix}.
\]
It satisfies $W^2=I$, $\det W=1$, $WAW=A^{-1}$,
$WBW=B^{-1}$ and $W^tJW=-J$. These identities give the actual
integral comparison of the two companion conventions, including
the pairing sign. For the based cusp used in the statement take
$c'=T_1^{-1}cT_1=(T_2T_1)^{-1}$. Then
\[
 \mu c\mu^{-1}=B^{-1}A,\qquad
 \mu c'\mu^{-1}=A(B^{-1}A)A^{-1}=AB^{-1}.
\]
The homology matrices are the inverse transposes under the actual
evaluation pairing with $L$.

\item Apply $\bigwedge^2\mu$ to the integral wedge lattice. In the
lexicographic basis of \cref{night:P1.2}, the invariant vector is
$\omega=(1,1,1,1,1,1)^t$. Let $Q$ have first column $\omega$ and
remaining columns the last five standard vectors, as in that result.
The induced quotient is the actual map
\[
 [u\wedge v]\longmapsto
 \operatorname{pr}_{2,\ldots,6}
 Q^{-1}\bigl(\mu(u)\wedge\mu(v)\bigr).
\]
Its kernel before passage to the quotient is precisely the invariant
line. Its inverse is induced by $(\bigwedge^2\mu^{-1})Q$ on the last
five standard vectors. The finite quotient matrices are the lower
five-by-five blocks of
\[
 Q^{-1}(\bigwedge^2 A^{-1})Q,\qquad
 Q^{-1}(\bigwedge^2 B)Q.
\]
They are exactly the integral quotient matrices of
\cref{night:P1.2}, and the cusp follows by the same based word
$(T_2T_1)^{-1}$. The invariant forms, integral inverses, peripheral
words and exterior minors are independently reproduced in
\texttt{P1\_7\_integral\_marking\_fresh.py} in the directory 
\texttt{code/continuation\_root/}.
\end{enumerate}
\end{proof}
